\documentclass[10pt]{article}

% ---------- IDIOMA Y CODIFICACIÓN ----------
\usepackage[utf8]{inputenc}
\usepackage[spanish]{babel}
\usepackage{amsmath}
\usepackage[margin=2.5cm]{geometry}
\usepackage{setspace}
\usepackage{graphicx}
\usepackage{listings}
\usepackage{minted}
\usepackage{tikz}
\usetikzlibrary{circuits.logic.US,babel}

% ---------- TABLAS (Actividad 1: identificación de dispositivos) ----------
\usepackage{booktabs}   % líneas de tabla profesionales (\toprule, \midrule, \bottomrule)
\usepackage{array}
\usepackage{ragged2e}
\usepackage[table]{xcolor} % opcional, para color en encabezado
\usepackage{tabularx}   % columnas que se ajustan al ancho de página
\usepackage{longtable}  % tablas que pueden partirse entre páginas (útil si tu tabla de dispositivos es larga)
\usepackage{caption}    % control de estilo de "Tabla X: ..."
\usepackage{float}
\usepackage{makecell}

% ---------- GLOSARIO (Actividad 2) ----------
% Enfoque simple con "description", sin necesidad de correr makeglossaries aparte
% (evita agregar un paso extra a tu pipeline de compilación).
\usepackage{enumitem}
\newlist{glosario}{description}{1}
\setlist[glosario]{
    font=\normalfont\bfseries,
    leftmargin=0pt,
    labelsep=0.5em,
    itemsep=0.5em
}

% ---------- Definiciones auxiliares ----------
\lstdefinestyle{asm}{
  language={[x86masm]Assembler},
  basicstyle=\ttfamily\small,
  keywordstyle=\bfseries,
  commentstyle=\itshape,
  numbers=left,
  numberstyle=\tiny,
  numbersep=6pt,
  frame=single,
  breaklines=true,
  showstringspaces=false,
  tabsize=4
}
%---- MINTED..----- 

\usepackage{xcolor}
\usepackage{minted}
\definecolor{bg}{rgb}{0.90,0.91,0.83}

\providecommand{\espacioimagen}[1]{%
  \fbox{\parbox[c][4.5cm][c]{0.9\linewidth}{\centering\itshape Insertar captura: #1}}}
 
% \figuraimagen{archivo.png}{Leyenda}{etiqueta}{descripción del recuadro}
\providecommand{\figuraimagen}[4]{%
  \begin{figure}[htbp]
    \centering
    \IfFileExists{img/#1}{\includegraphics[width=0.9\linewidth]{img/#1}}{\espacioimagen{#4\\[2pt]{\upshape\small\texttt{img/#1}}}}
    \caption{#2}\label{#3}
  \end{figure}}
 

%      \relleno{Tu texto}    -> muestra tu texto en rojo (útil para ir dejando borradores)
\newcommand{\relleno}[1][\textbf{[COMPLETAR AQUÍ]}]{\textcolor{red}{#1}}
% ---------- CONFIGURACIONES PARA ENCABEZADO Y PORTADA ----------
\usepackage{colortbl}
\usepackage[colorlinks=true,linkcolor=black,urlcolor=black,citecolor=black]{hyperref}

% Definición de colores y columnas personalizadas (usadas en la portada)
\definecolor{tablebackground}{rgb}{0.8, 0, 0}
\newcolumntype{x}[1]{>{\centering\arraybackslash\hspace{0pt}}p{#1}}
\newcolumntype{M}[1]{>{\centering\arraybackslash}m{#1}}
\newcolumntype{N}{@{}m{0pt}@{}}
\newcolumntype{P}[1]{>{\RaggedRight\arraybackslash}p{#1}}

% Variables de datos para la portada y encabezado
\newcommand{\itemEmail}{ccapatinta@unsa.edu.pe}
\newcommand{\itemStudent}{Camilo B.A. Capatinta Almanza (Teoría grupo B)}
\newcommand{\itemCourse}{Lab. Arq. de computadoras}
\newcommand{\itemCourseCode}{1702227}
\newcommand{\itemSemester}{IV}
\newcommand{\itemUniversity}{Universidad Nacional de San Agustín de Arequipa}
\newcommand{\itemFaculty}{Facultad de Ingeniería de Producción y Servicios}
\newcommand{\itemDepartment}{Departamento Académico de Ingeniería de Sistemas e Informática}
\newcommand{\itemSchool}{Escuela Profesional de Ingeniería de Sistemas}
\newcommand{\itemAcademic}{2026 - B}
\newcommand{\itemInput}{20/10/2026}
\newcommand{\itemOutput}{--/10/2026}
\newcommand{\itemPracticeNumber}{05}
\newcommand{\itemTheme}{Programación para el microprocesador INTEL - Instrucciones Aritméticas}

% Configuración del encabezado y pie de página (fancyhdr)
\usepackage{fancyhdr}
\pagestyle{fancy}
\fancyhf{}
\setlength{\headheight}{40.51407pt}
\renewcommand{\headrulewidth}{1pt}
\renewcommand{\footrulewidth}{1pt}

% Encabezado
\fancyhead[L]{\raisebox{-0.2\height}{\includegraphics[width=3cm]{img/logo_episunsa.png}}}
\fancyhead[C]{\fontsize{7}{7}\selectfont    \itemUniversity \\ \itemFaculty \\ \itemDepartment \\ \itemSchool \\ \textbf{\itemCourse}}
\fancyhead[R]{\raisebox{-0.2\height}{\includegraphics[width=1.2cm]{img/logo_abet.png}}}

% Pie de página
\fancyfoot[L]{Estudiante Camilo Capatinta}
\fancyfoot[C]{\itemCourse}
\fancyfoot[R]{Página \thepage}

\begin{document}

\section{Actividades}

\subsection{Parte 1: Verdadero (V) o Falso (F)}

\begin{itemize}
  \item \textbf{1. Falso (F).} La SRAM es muy rápida, pero es costosa y de baja densidad; por eso se usa como memoria \emph{caché}. La memoria principal de un computador se construye con DRAM, que es más barata y densa.

  \item \textbf{2. Verdadero (V).} Los procesadores RISC son muy rápidos gracias a su juego reducido de instrucciones, pero necesitan del orden de un 30\,\% más de instrucciones que una máquina CISC para un programa tipo.

  \item \textbf{3. Verdadero (V).} La memoria virtual requiere traducir direcciones lógicas a físicas, tarea que realiza la MMU (\emph{Memory Management Unit}), ya sea integrada en el procesador o como dispositivo aparte.

  \item \textbf{4. Verdadero (V).} Describe correctamente el funcionamiento de la caché: se transfiere un bloque en el primer acceso, los accesos siguientes al bloque se leen de la caché, la función de mapeo define la correspondencia bloque--memoria principal y el algoritmo de reemplazo decide qué bloque se elimina cuando la caché se llena.

  \item \textbf{5. Verdadero (V).} En una caché de correspondencia directa, los bits centrales de la dirección (índice) seleccionan el bloque, la etiqueta se compara con la almacenada y, si coinciden y el bit de válido está activo, hay acierto; en caso contrario hay fallo y se accede a la memoria principal.

  \item \textbf{6. Falso (F).} El orden dado (Componente, Electrónico, Digital, RTL, PMS) corresponde al modelo \emph{bottom-up}. El orden \emph{top-down} correcto es: PMS, RTL, Digital, Electrónico y Componente.

  \item \textbf{7. Falso (F).} Las memorias Flash son de lectura/escritura y de acceso aleatorio, pero son \textbf{no volátiles}: conservan la información sin alimentación.

  \item \textbf{8. Verdadero (V).} La paginación y la segmentación permiten un espacio de direccionamiento virtual mayor que la memoria física, no requieren tener todo el programa en memoria y facilitan la compartición de código y los programas reubicables.

  \item \textbf{9. Falso (F).} El registro de instrucción (IR) almacena la \emph{instrucción actual} que se está ejecutando. La dirección de la siguiente instrucción la almacena el contador de programa (PC).

  \item \textbf{10. Verdadero (V).} Un sistema operativo es el conjunto de programas, instalados en software o firmware, que hace utilizable el hardware.

  \item \textbf{11. Falso (F).} Un computador es un sistema secuencial \textbf{síncrono} (sincronizado por una señal de reloj), no asíncrono. El resto del enunciado (representación binaria, 1 lógico $\approx$ 5\,V y 0 lógico $\approx$ 0\,V) es correcto.

  \item \textbf{12. Verdadero (V).} El complemento a dos de $0Ah$ se obtiene así:
  \begin{align*}
    0Ah &= 0000\ 1010_2 \\
    \text{complemento a uno} &= 1111\ 0101_2 \\
    +1 &\;\Rightarrow\; 1111\ 0110_2
  \end{align*}
  que coincide con el valor del enunciado ($F6h = -10_{10}$).
\end{itemize}

\subsection{Parte 1: Preguntas de desarrollo}

\begin{itemize}
  \item \textbf{13. Capacidad de direccionamiento: 16 KBytes.}

  Dirección inicial: $04000h$; dirección final: $07FFFh$. Como cada dirección contiene un byte:
  \[
    07FFFh - 04000h + 1 = 03FFFh + 1 = 4000h = 2^{14}\ \text{bytes}
  \]
  \[
    2^{14}\ \text{bytes} = 2^{4} \cdot 2^{10} = 16\ \text{KBytes}
  \]

  \item \textbf{14. Número de bits del bus de direcciones: 39 bits.}

  Si cada dirección contiene un byte, hay que direccionar $512\ \text{GBytes}$:
  \[
    512\ \text{GB} = 2^{9} \cdot 2^{30}\ \text{bytes} = 2^{39}\ \text{bytes}
  \]
  Por tanto, se necesitan $n = 39$ bits en el bus de direcciones.

  \item \textbf{15. Resultado: $AC = 0$ y contenido de la dirección $401 = 10$.}

  Estado inicial: $M[400] = 10$, $M[401] = 5$. El programa funciona como un bucle: carga $M[401]$, lo incrementa, lo guarda de nuevo y resta $M[400]$; si el resultado es cero, termina (\texttt{JMZ fin}); si no, repite (\texttt{JMP bucle}).

  \begin{center}
  \begin{tabular}{c c c c c}
    \hline
    \textbf{Iteración} & \textbf{LW \#401} & \textbf{INC AC / SW \#401} & \textbf{SUB \#400} & \textbf{¿Salta a fin?} \\
    \hline
    1  & $AC=5$  & $AC=6,\ M[401]=6$   & $AC=6-10=-4$ & No \\
    2  & $AC=6$  & $AC=7,\ M[401]=7$   & $AC=-3$      & No \\
    3  & $AC=7$  & $AC=8,\ M[401]=8$   & $AC=-2$      & No \\
    4  & $AC=8$  & $AC=9,\ M[401]=9$   & $AC=-1$      & No \\
    5  & $AC=9$  & $AC=10,\ M[401]=10$ & $AC=0$       & Sí \\
    \hline
  \end{tabular}
  \end{center}

  Al llegar a \texttt{fin: STP}, el acumulador vale $AC = 0$ (resultado de $10 - 10$) y la posición $401$ contiene $10$.
\end{itemize}

\end{document}